% Part 2 should therefore focus on static approaches to representing accessibility or evolutionary potential. This can include classic quantitative genetics (e.g. the G matrix), with co-expression analyses treated as illustrative examples rather than the main focus. Here we can highlight both the strengths of these approaches and their limitations, such as the fact that many different underlying topologies can map to the same covariation structure.

- What evolutionary processes determine accessibility in phenotypic landscapes? 
- What determines evolutionary trajectories? Long term, short term
- Dynamics of adaptive landscapes, do the optima shift? Which processes lead to this?
    - Frequency-dependent selection
    - Coevolution
    - Environmental shifts
- Standing genetic variation can shape short-term response; therefore, the G-matrix structure is central.
- Variability (generation of new variation) determines the structure of G. New variation can be shaped by: distribution of new mutations (mutational targets, canalization), development (not constraints, just development), internal selection.
- Final covariance pattern is degenerate (many to one), what are the consequences of this for short and long-term prediction and inference?
- The stability of the variability generating mechanisms becomes a central question. G-matrices can change very fast, but don't appear to. What are the explanations for this? How can we test them? Mechanistic analysis might be the key ingredient.
- Recovering GRN from covariation data, regulatory relations should be more stable than standing variation. Mechanism over realized variation.

% old intro

Evolution by natural selection occurs via the differential survival and reproduction of heritable phenotypic variation. 
Both the supply of phenotypic variation and its selection are therefore integral parts of attempts to model and predict evolution.
However, the developmental mechanisms that define the genotype-phenotype mapping and the organism-environment interactions that determine fitness can be conceptually intricate and empirically difficult to determine. 
Research programs in evolutionary biology have tended to emphasize one over the other and employ strategic simplifications to manage the otherwise unwieldy complexity of real-world evolution. 
On the one hand, the ‘phenotypic gambit’ is widely used in optimality approaches and evolutionary game theory, assuming that the underlying genetic architectures and developmental mechanisms in the long run give way to natural selection and can therefore be neglected in predicting evolutionary outcomes and adaptation. 
On the other hand, models in evo-devo and self-organization are primarily concerned with the mechanisms that define the genotype-phenotype map, carrying the implicit assumption that phenotypes evolve primarily in ways driven by the underlying mechanisms, not because of their selective benefit. 

While both approaches represent legitimate and complementary perspectives on evolutionary biology, there exists a perennial tension between the two whether to interpret evolutionary outcomes as primarily driven by the creative force of natural selection, or the potentialities afforded by genetic and developmental mechanisms.
How far the conclusions of the one approach can be pushed without incorporating the other is an open question addressed by largely separate research communities. 

One field that contains this conflict as a whole is quantitative genetics, a field that started with the breeders’ science of quantitative traits, but has since evolved into a broad-purpose toolkit applied in phylogenetic inference, genome-wide association studies and eco-evolutionary dynamics.
A core concept of the field is the G-matrix, which captures the additive variance and covariance of a suite of quantitative traits present in a population.
In many cases, the G-matrix has been shown to be remarkably constant, even over large time frames and across distantly related taxa. Other studies however have shown G to vary strongly even over a few generations.
The debate around the constancy of the G-matrix reflects the conflict between a focus on adaptation and a focus on underlying mechanisms discussed earlier.
Under conditions where the G-matrix is constant, the G-matrix basically captures the properties of the genotype-phenotype map to a sufficient degree.
Where it is not constant, we need to dig deeper into the underlying mechanisms of development to be able to model the supply of phenotypic variation to the process of natural selection. Several models have occupied themselves with this question, studying the impact of epistasis and pleiotropy \parencite{gjuvsland2007statistical, hether2014genetic, jiang2025alternative}, etc. 

